\subsection*{状态不再是几个数，而是一整条曲线}
一根杆在各位置有不同温度，所以状态是温度分布 $u(x,t)$。时间改变时，整条曲线都变。偏导数让我们分别固定位置观察时间变化，或固定时间观察空间变化，两者承担不同物理角色。

为什么热方程出现空间二阶导数，而不是温度本身？一个位置比两侧邻居都热，就会向两侧散热；比邻居冷则吸热。决定变化的是中心相对于邻域平均的偏离。三点模型给出 $u(x-h)-2u(x)+u(x+h)$，除以 $h^2$ 后趋向曲率。

\subsection*{空间斜率与温度变化之间隔着一层差分}
在常系数模型中，热通量与负温度梯度成正比；某点温度变化来自流入与流出的差，所以对通量再取一次空间变化，便得到二阶导数。线性温度分布虽然有热流，但各小段流入流出平衡，内部温度可以暂时不变。

端点怎样处理会改变整个问题。保持端点温度固定，与隔热端点无热流，是不同边界条件；同一个初始曲线在不同边界下会演化不同。只给方程而不给初始和边界信息，通常不足以选出具体解。

若做离散模拟，还要检查时间步稳定性。把空间分得更细却不相应调整时间步，可能让数值解振荡甚至爆炸，那不一定是物理现象。接下来我们寻找一些形状简单、能精确追踪的模式，用来解释也用来校验模拟。
